\documentclass[10pt,a4paper]{amsart}
\usepackage{amsmath,amssymb,amsthm,graphicx}
\usepackage{verbatim}
\newenvironment{reference}{\comment}{\endcomment}
%\newenvironment{reference}{}{}
\newenvironment{slogan}{\comment}{\endcomment}
\newenvironment{history}{\comment}{\endcomment}

% For commutative diagrams we use Xy-pic
\usepackage[all]{xy}

% We use 2cell for 2-commutative diagrams.
\xyoption{2cell}
% Theorem environments.
%
\theoremstyle{plain}
\newtheorem{theorem}[subsection]{Theorem}
\newtheorem{proposition}[subsection]{Proposition}
\newtheorem{lemma}[subsection]{Lemma}

\theoremstyle{definition}
\newtheorem{definition}[subsection]{Definition}
\newtheorem{example}[subsection]{Example}
\newtheorem{exercise}[subsection]{Exercise}
\newtheorem{situation}[subsection]{Situation}

\theoremstyle{remark}
\newtheorem{remark}[subsection]{Remark}
\newtheorem{remarks}[subsection]{Remarks}

\numberwithin{equation}{subsection}

% Macros
%
\def\lim{\mathop{\mathrm{lim}}\nolimits}
\def\colim{\mathop{\mathrm{colim}}\nolimits}
\def\Spec{\mathop{\mathrm{Spec}}}
\def\Hom{\mathop{\mathrm{Hom}}\nolimits}
\def\Ext{\mathop{\mathrm{Ext}}\nolimits}
\def\SheafHom{\mathop{\mathcal{H}\!\mathit{om}}\nolimits}
\def\SheafExt{\mathop{\mathcal{E}\!\mathit{xt}}\nolimits}
\def\Sch{\mathit{Sch}}
\def\Mor{\mathop{\mathrm{Mor}}\nolimits}
\def\Ob{\mathop{\mathrm{Ob}}\nolimits}
\def\Sh{\mathop{\mathit{Sh}}\nolimits}
\def\NL{\mathop{N\!L}\nolimits}
\def\CH{\mathop{\mathrm{CH}}\nolimits}
\def\proetale{{pro\text{-}\acute{e}tale}}
\def\etale{{\acute{e}tale}}
\def\QCoh{\mathit{QCoh}}
\def\Ker{\mathop{\mathrm{Ker}}}
\def\Im{\mathop{\mathrm{Im}}}
\def\Coker{\mathop{\mathrm{Coker}}}
\def\Coim{\mathop{\mathrm{Coim}}}

% Boxtimes
%
\DeclareMathSymbol{\boxtimes}{\mathbin}{AMSa}{"02}

%
% Macros for moduli stacks/spaces
%
\def\QCohstack{\mathcal{QC}\!\mathit{oh}}
\def\Cohstack{\mathcal{C}\!\mathit{oh}}
\def\Spacesstack{\mathcal{S}\!\mathit{paces}}
\def\Quotfunctor{\mathrm{Quot}}
\def\Hilbfunctor{\mathrm{Hilb}}
\def\Curvesstack{\mathcal{C}\!\mathit{urves}}
\def\Polarizedstack{\mathcal{P}\!\mathit{olarized}}
\def\Complexesstack{\mathcal{C}\!\mathit{omplexes}}
% \Pic is the operator that assigns to X its picard group, usage \Pic(X)
% \Picardstack_{X/B} denotes the Picard stack of X over B
% \Picardfunctor_{X/B} denotes the Picard functor of X over B
\def\Pic{\mathop{\mathrm{Pic}}\nolimits}
\def\Picardstack{\mathcal{P}\!\mathit{ic}}
\def\Picardfunctor{\mathrm{Pic}}
\def\Deformationcategory{\mathcal{D}\!\mathit{ef}}
\usepackage{cleveref}
\crefname{originalreference}{original reference}{original references}
\Crefname{originalreference}{Original reference}{Original references}
\makeatletter
\newlabel{lemma-normal-closure}{{09DT}{source}{}{sourcepointer-1}{}}
\newlabel{lemma-normal-closure@cref}{{[originalreference][0][]09DT}{[1][0][]source}{}{}{}}
\makeatother
\crefname{originalreference}{original reference}{original references}
\Crefname{originalreference}{Original reference}{Original references}
\makeatletter
\newlabel{lemma-normally-generated}{{0BR3}{source}{}{sourcepointer-2}{}}
\newlabel{lemma-normally-generated@cref}{{[originalreference][0][]0BR3}{[1][0][]source}{}{}{}}
\makeatother
\title{Inside a normal algebraic extension, finite subextensions have finite normal closures, and normality survives finite steps}
\author{The Stacks Project Authors}
\date{stacks-project:fields.tex:local-snapshot-20261012; GNU FDL 1.2 or later}
\begin{document}
\maketitle
\noindent\textit{Editorial note.} All delivered mathematical text is the author's own, copied from the pinned Stacks Project chapter file fields.tex; no statement or proof step was rewritten, reordered or removed.\par\medskip
\noindent\textit{Editorial note.} The counted proof is the author's complete argument: Item (2) needs M/K normal (not just M'/K), which is exactly what the hypothesis provides; this lemma feeds the Galois-uniqueness machinery for normal closures inside normal towers.\par\medskip
\noindent\textit{Editorial note.} The label lemma-normal-closure-inside-normal names the counted statement, and the proof environment immediately following it in the source is delivered verbatim as its proof.\par\medskip
\noindent\textit{Editorial note.} The counted statement is carried by the author's own lemma environment in the chapter source.\par\medskip
\noindent\textit{Allowed standard prerequisites (editorial).}\begin{itemize}
\item Same chapter: lemma-normal-closure (delivered card L48334: the normal closure of a finite extension is finite and normal), lemma-normally-generated (delivered card L48362 in this batch: adjoining generators whose minimal polynomials split completely gives a normal extension), L/K normal makes every irreducible over K with a root in L split in L
\end{itemize}\medskip
\ifdefined\hypertarget\hypertarget{sourcepointer-1}{}\fi
\noindent\textit{Source reference.} Original reference 09DT: sources/stacks-fields/pinned-original-source.tex, source lines 2303--2303. Prerequisite pointer invoked at 2364; it names a standard fact of the same chapter that the counted argument uses, and every step of the counted proof built on it is retained verbatim in the delivered body.\par\medskip
\ifdefined\hypertarget\hypertarget{sourcepointer-2}{}\fi
\noindent\textit{Source reference.} Original reference 0BR3: sources/stacks-fields/pinned-original-source.tex, source lines 2139--2139. Prerequisite pointer invoked at 2372; it names a standard fact of the same chapter that the counted argument uses, and every step of the counted proof built on it is retained verbatim in the delivered body.\par\medskip
\begin{lemma}
\label{lemma-normal-closure-inside-normal}
Let $L/K$ be an algebraic normal extension.
\begin{enumerate}
\item If $L/M/K$ is a subextension with $M/K$ finite, then there exists
a tower $L/M'/M/K$ with $M'/K$ finite and normal.
\item If $L/M'/M/K$ is a tower with $M/K$ normal and $M'/M$ finite,
then there exists a tower $L/M''/M'/M/K$ with $M''/M$
finite and $M''/K$ normal.
\end{enumerate}
\end{lemma}

\begin{proof}
Proof of (1). Let $M'$ be the smallest subextension of $L/K$ containing $M$
which is normal over $K$. By Lemma \ref{lemma-normal-closure}
this is the normal closure of $M/K$ and is finite over $K$.

\medskip\noindent
Proof of (2). Let $\alpha_1, \ldots, \alpha_n \in M'$ generate $M'$ over $M$.
Let $P_1, \ldots, P_n$ be the minimal polynomials of
$\alpha_1, \ldots, \alpha_n$ over $K$. Let $\alpha_{i, j}$ be the roots
of $P_i$ in $L$. Let $M''  = M(\alpha_{i, j})$. It follows from
Lemma \ref{lemma-normally-generated}
(applied with the set of generators $M \cup \{\alpha_{i, j}\}$)
that $M''$ is normal over $K$.
\end{proof}

\end{document}
